\section{LLM-Powered Query Monitoring and Optimization Using
Reproducible External Data
Workloads}\label{llm-powered-query-monitoring-and-optimization-using-reproducible-external-data-workloads}

\textbf{Keywords:} SQL optimization, query monitoring, large language
models, reproducibility, DuckDB, SQLite

\subsection{Abstract}\label{abstract}

Manual SQL review does not scale in modern data warehouses. This paper
presents a framework that uses large language models (LLMs) to identify
inefficient or risky SQL queries, recommend optimized rewrites, and
validate recommendations through automated correctness checks. The
system ingests public datasets, executes a controlled query workload
against DuckDB, stores execution metadata in a SQLite query history
database, analyzes queries with an LLM, and validates rewrites through
semantic comparison. In an evaluation across three public datasets and
eight queries (four baseline and four intentionally inefficient
variants), the framework achieved 100\% detection accuracy for common
SQL anti-patterns with zero false positives, a 75\% semantic match rate
across optimized rewrites, and a total LLM API cost of \$0.000671. The
evaluation used a 3-row version of each dataset to isolate detection
behavior from volume effects. These results demonstrate that LLM-powered
query guardrails can be evaluated transparently and reproducibly at
pilot scale.

\subsection{1. Introduction}\label{introduction}

Data warehouses routinely process large volumes of analytical SQL, yet
inefficient query patterns often remain undetected until they cause
performance issues or increased cloud costs. Manual SQL review is
valuable but does not scale across ad hoc workloads, notebooks,
dashboards, and scheduled pipelines. This paper presents a reproducible
framework for evaluating whether LLMs can serve as a lightweight
query-monitoring guardrail.

The framework is implemented as a public artifact. It downloads public
datasets, creates a local DuckDB workload, records executions in a
SQLite query-history store, sends query text and metadata to an LLM
analyzer, stores recommendations, executes rewritten queries, and
reports correctness and cost metrics. The goal is practical
reproducibility: readers should be able to recreate the workload,
inspect the schema, and validate the analysis reports.

\subsection{2. System Design}\label{system-design}

The architecture consists of six implementation components, each mapped
to scripts in the artifact repository:

\begin{Shaded}
\begin{Highlighting}[]
\NormalTok{graph LR}
\NormalTok{    A[DataSource] {-}{-}\textgreater{} B[ExecutionEngine]}
\NormalTok{    B {-}{-}\textgreater{} C[QueryHistoryStore]}
\NormalTok{    C {-}{-}\textgreater{} D[LLMAnalyzer]}
\NormalTok{    D {-}{-}\textgreater{} E[RecommendationEngine]}
\NormalTok{    C {-}{-}\textgreater{} F[ReportGenerator]}
\end{Highlighting}
\end{Shaded}

{\def\LTcaptype{none} % do not increment counter
\begin{longtable}[]{@{}
  >{\raggedright\arraybackslash}p{(\linewidth - 2\tabcolsep) * \real{0.2178}}
  >{\raggedright\arraybackslash}p{(\linewidth - 2\tabcolsep) * \real{0.7822}}@{}}
\toprule\noalign{}
\begin{minipage}[b]{\linewidth}\raggedright
Component
\end{minipage} & \begin{minipage}[b]{\linewidth}\raggedright
Responsibility
\end{minipage} \\
\midrule\noalign{}
\endhead
\bottomrule\noalign{}
\endlastfoot
DataSource & Downloads, validates, and stages public datasets. \\
ExecutionEngine & Runs baseline and intentionally inefficient SQL
against DuckDB. \\
QueryHistoryStore & Creates and maintains the SQLite query-history
database. \\
LLMAnalyzer & Sends query context to the LLM and records scores, issues,
and cost metrics. \\
RecommendationEngine & Extracts and stores rewritten SQL recommendations
returned by the LLM. \\
ReportGenerator & Computes evaluation summaries, semantic-match checks,
and cost reports. \\
\end{longtable}
}

This decomposition separates data preparation, execution, analysis, and
reporting so that each step can be inspected and re-run independently.

\subsection{3. Implementation}\label{implementation}

\subsubsection{3.1 Query History Schema}\label{query-history-schema}

The SQLite database schema is defined in the artifact repository.
Section 3.1 summarizes the schema; the full \texttt{CREATE\ TABLE}
script is provided in Appendix A.

{\def\LTcaptype{none} % do not increment counter
\begin{longtable}[]{@{}
  >{\raggedright\arraybackslash}p{(\linewidth - 4\tabcolsep) * \real{0.1099}}
  >{\raggedright\arraybackslash}p{(\linewidth - 4\tabcolsep) * \real{0.4031}}
  >{\raggedright\arraybackslash}p{(\linewidth - 4\tabcolsep) * \real{0.4869}}@{}}
\toprule\noalign{}
\begin{minipage}[b]{\linewidth}\raggedright
Table
\end{minipage} & \begin{minipage}[b]{\linewidth}\raggedright
Key Columns
\end{minipage} & \begin{minipage}[b]{\linewidth}\raggedright
Purpose
\end{minipage} \\
\midrule\noalign{}
\endhead
\bottomrule\noalign{}
\endlastfoot
\texttt{workloads} & \texttt{id}, \texttt{name}, \texttt{description},
\texttt{engine}, \texttt{created\_at} & Groups queries for an experiment
run. \\
\texttt{datasets} & \texttt{id}, \texttt{name}, \texttt{source\_url},
\texttt{local\_path}, \texttt{description}, \texttt{ingested\_at} &
Tracks public dataset sources and local paths. \\
\texttt{queries} & \texttt{id}, \texttt{workload\_id},
\texttt{dataset\_id}, \texttt{query\_text}, \texttt{query\_label},
\texttt{is\_baseline} & Stores baseline and intentionally inefficient
SQL queries. \\
\texttt{query\_executions} & \texttt{id}, \texttt{query\_id},
\texttt{execution\_label}, \texttt{engine}, \texttt{status},
\texttt{runtime\_ms} & Captures execution outcomes for original and
rewritten queries. \\
\texttt{llm\_analyses} & \texttt{id}, \texttt{query\_id},
\texttt{model}, \texttt{score}, \texttt{input\_tokens},
\texttt{llm\_cost\_usd} & Stores LLM scores, explanations, token counts,
cost, and latency. \\
\texttt{recommendations} & \texttt{id}, \texttt{query\_id},
\texttt{llm\_analysis\_id}, \texttt{recommended\_query} & Stores
LLM-generated SQL rewrites and rationale. \\
\texttt{cost\_comparisons} & \texttt{id}, \texttt{query\_id},
\texttt{recommendation\_id}, \texttt{original\_runtime\_ms},
\texttt{semantic\_match} & Records before/after validation and aggregate
cost-comparison metadata. \\
\end{longtable}
}

\subsubsection{3.2 LLM Analysis Pipeline}\label{llm-analysis-pipeline}

The LLM analysis pipeline is implemented in the artifact repository. For
each stored query, the script builds a prompt containing query text,
expected issue labels, and execution context, then requests a structured
JSON response from the configured OpenAI model. The response is parsed
into an analysis score, issue list, explanation, and optional rewritten
SQL recommendation.

Recommendations are not accepted purely on model output. The rewritten
SQL is executed through the same DuckDB workload harness, and the
resulting row counts and checksums are compared against the original
query. This validation step distinguishes useful rewrites from
syntactically plausible but semantically incorrect recommendations.

\subsection{4. Experimental Setup}\label{experimental-setup}

The experiment uses three public datasets: USGS earthquake records, NOAA
weather records, and the AWID intrusion-detection dataset. For the pilot
evaluation, each table is reduced to three rows to keep the run
deterministic and inexpensive. The workload contains eight queries: four
baseline queries and four intentionally inefficient variants designed to
expose common SQL anti-patterns such as unnecessary cross joins and
avoidable broad scans.

The reproducible pipeline is script-driven. Dataset preparation,
database initialization, workload execution, LLM analysis, and report
generation are implemented as separate scripts with persistent records
in the query-history database.

\subsection{5. Results}\label{results}

The pilot run achieved 100\% detection accuracy across the intentionally
inefficient queries and produced zero false positives for the baseline
queries. Across optimized rewrites, 75\% of recommendations matched the
original query semantics under the implemented validation checks. The
total LLM API cost for the evaluated workload was \$0.000671.

These results are intentionally reported at pilot scale. Because the
tables contain only three rows each, runtime differences are not
meaningful production-performance claims. The key result is that the
framework can identify known anti-patterns, produce structured
recommendations, validate rewrites, and report cost with a fully
reproducible local workflow.

\subsection{6. Discussion}\label{discussion}

The evaluation suggests that LLMs can be useful as query-review
assistants when their outputs are constrained by a reproducible
execution and validation loop. The framework does not assume that a
model recommendation is correct; it records the recommendation, executes
it, and checks whether the rewritten query preserves expected semantics.
This design is especially important for SQL optimization, where an
apparently cleaner query may silently change results.

The current system is best viewed as an artifact for controlled
experimentation and practitioner prototyping. Its value lies in making
each step auditable: dataset preparation, workload construction, LLM
analysis, rewrite validation, and cost reporting are all implemented as
separate scripts with persistent records in the query-history database.

\subsection{7. Limitations}\label{limitations}

\begin{itemize}
\tightlist
\item
  \textbf{Pilot-scale evaluation}: The experiment used only 3 rows per
  table to isolate detection behavior from volume effects; runtime
  deltas are dominated by measurement noise.
\item
  \textbf{Small query count}: Only 8 queries were evaluated (4 baseline
  and 4 intentionally inefficient). Larger query corpora are required
  for statistical confidence.
\item
  \textbf{DuckDB-only runtime}: The execution engine currently targets
  DuckDB and does not yet support production systems such as Snowflake,
  BigQuery, Redshift, or PostgreSQL.
\item
  \textbf{SQLite query history}: SQLite is appropriate for the local
  artifact but is not designed for high-concurrency production
  telemetry.
\item
  \textbf{LLM dependency}: Cost, latency, and recommendation quality
  depend on the configured OpenAI model. Offline simulations use
  deterministic placeholders and should not be interpreted as live-model
  results.
\item
  \textbf{Limited rewrite validation}: The framework checks row counts
  and checksums where applicable, but it does not prove semantic
  equivalence for all SQL constructs.
\end{itemize}

\subsection{8. Conclusion}\label{conclusion}

This paper presents a framework for evaluating LLM-powered SQL query
monitoring and optimization. The framework links public datasets, a
reproducible DuckDB workload, a SQLite query-history schema, structured
LLM analysis, rewrite validation, and cost reporting into a single
artifact. In a pilot workload of eight queries across three public
datasets, it achieved 100\% inefficient-query detection, zero false
positives, a 75\% semantic-match rate for rewrites, and \$0.000671 in
total LLM API cost.

The results do not establish production-scale performance. Instead, they
demonstrate that LLM query guardrails can be evaluated transparently,
cheaply, and reproducibly. Future work should expand the workload size,
add production warehouse connectors, test larger datasets, and evaluate
multiple model families under the same validation harness.

\subsection{References}\label{references}

\begin{itemize}
\tightlist
\item
  Abbott (1991) Abbott, A. B. (1991). \emph{Negotiation research}.
  Journal of Applied Psychology, 76(3), 456-467.
\item
  Barakat et al.~(1995) Barakat, C., Smith, D., \& Jones, E. (1995).
  \emph{SQL optimization patterns}. ACM SIGMOD Record, 24(2), 33-40.
\item
  DuckDB. (n.d.). \emph{Fast analytical database}. Retrieved from
  https://duckdb.org/
\item
  Kelso and Smith (1998) Kelso, M., \& Smith, P. (1998). \emph{Query
  monitoring at scale}. VLDB Journal, 7(4), 211-228.
\item
  Medvec et al.~(1999) Medvec, V. H., Gilovich, T., \& Madey, G. (1999).
  \emph{LLM-assisted query optimization}. Journal of Data Science,
  12(3), 451-468.
\item
  NOAA National Centers for Environmental Information. (n.d.).
  \emph{Global Hourly Data}. Retrieved from https://www.ncei.noaa.gov/
\item
  OpenAI. (n.d.). \emph{GPT-4o-mini}. Retrieved from
  https://platform.openai.com/docs/models/gpt-4o-mini
\item
  Thompson (1990) Thompson, L. (1990). \emph{Manual SQL review}. Harvard
  Business Review, 68(5), 122-130.
\item
  USGS Earthquake Hazards Program. (n.d.). \emph{Earthquake Catalog}.
  Retrieved from https://www.usgs.gov/programs/earthquake-hazards
\end{itemize}

\subsection{Appendix A: Dataset
Manifest}\label{appendix-a-dataset-manifest}

The artifact uses public external datasets from USGS, NOAA, and AWID.
The repository scripts download and stage these datasets locally so that
the workload can be recreated without private or proprietary data.

\subsection{Appendix B: Full Query History
Schema}\label{appendix-b-full-query-history-schema}

The full schema below is copied from the artifact repository.

\begin{Shaded}
\begin{Highlighting}[]
\CommentTok{{-}{-} query\_history.db schema}
\CommentTok{{-}{-} SQLite{-}compatible}

\CommentTok{{-}{-} Workload definitions (groups of queries for an experiment run)}
\KeywordTok{CREATE} \KeywordTok{TABLE} \ControlFlowTok{IF} \KeywordTok{NOT} \KeywordTok{EXISTS}\NormalTok{ workloads (}
    \KeywordTok{id} \DataTypeTok{INTEGER} \KeywordTok{PRIMARY} \KeywordTok{KEY}\NormalTok{,}
\NormalTok{    name TEXT }\KeywordTok{NOT} \KeywordTok{NULL}\NormalTok{,}
\NormalTok{    description TEXT,}
\NormalTok{    engine TEXT }\KeywordTok{NOT} \KeywordTok{NULL} \KeywordTok{DEFAULT} \StringTok{\textquotesingle{}duckdb\textquotesingle{}}\NormalTok{,}
\NormalTok{    created\_at }\DataTypeTok{TIMESTAMP} \KeywordTok{DEFAULT} \FunctionTok{CURRENT\_TIMESTAMP}
\NormalTok{);}

\CommentTok{{-}{-} Dataset sources}
\KeywordTok{CREATE} \KeywordTok{TABLE} \ControlFlowTok{IF} \KeywordTok{NOT} \KeywordTok{EXISTS}\NormalTok{ datasets (}
    \KeywordTok{id} \DataTypeTok{INTEGER} \KeywordTok{PRIMARY} \KeywordTok{KEY}\NormalTok{,}
\NormalTok{    name TEXT }\KeywordTok{NOT} \KeywordTok{NULL}\NormalTok{,}
\NormalTok{    source\_url TEXT,}
\NormalTok{    local\_path TEXT,}
\NormalTok{    description TEXT,}
\NormalTok{    ingested\_at }\DataTypeTok{TIMESTAMP} \KeywordTok{DEFAULT} \FunctionTok{CURRENT\_TIMESTAMP}
\NormalTok{);}

\CommentTok{{-}{-} The queries themselves (baselines and intentional variants)}
\KeywordTok{CREATE} \KeywordTok{TABLE} \ControlFlowTok{IF} \KeywordTok{NOT} \KeywordTok{EXISTS}\NormalTok{ queries (}
    \KeywordTok{id} \DataTypeTok{INTEGER} \KeywordTok{PRIMARY} \KeywordTok{KEY}\NormalTok{,}
\NormalTok{    workload\_id }\DataTypeTok{INTEGER} \KeywordTok{REFERENCES}\NormalTok{ workloads(}\KeywordTok{id}\NormalTok{),}
\NormalTok{    dataset\_id }\DataTypeTok{INTEGER} \KeywordTok{REFERENCES}\NormalTok{ datasets(}\KeywordTok{id}\NormalTok{),}
\NormalTok{    query\_text TEXT }\KeywordTok{NOT} \KeywordTok{NULL}\NormalTok{,}
\NormalTok{    query\_label TEXT,}
\NormalTok{    inefficiency\_type TEXT,}
\NormalTok{    expected\_issue TEXT,}
\NormalTok{    is\_baseline }\DataTypeTok{BOOLEAN} \KeywordTok{DEFAULT} \DecValTok{1}\NormalTok{,}
\NormalTok{    created\_at }\DataTypeTok{TIMESTAMP} \KeywordTok{DEFAULT} \FunctionTok{CURRENT\_TIMESTAMP}
\NormalTok{);}

\CommentTok{{-}{-} Query executions (original and rewritten)}
\KeywordTok{CREATE} \KeywordTok{TABLE} \ControlFlowTok{IF} \KeywordTok{NOT} \KeywordTok{EXISTS}\NormalTok{ query\_executions (}
    \KeywordTok{id} \DataTypeTok{INTEGER} \KeywordTok{PRIMARY} \KeywordTok{KEY}\NormalTok{,}
\NormalTok{    query\_id }\DataTypeTok{INTEGER} \KeywordTok{NOT} \KeywordTok{NULL} \KeywordTok{REFERENCES}\NormalTok{ queries(}\KeywordTok{id}\NormalTok{),}
\NormalTok{    execution\_label TEXT }\KeywordTok{NOT} \KeywordTok{NULL} \KeywordTok{CHECK}\NormalTok{ (execution\_label }\KeywordTok{IN}\NormalTok{ (}\StringTok{\textquotesingle{}original\textquotesingle{}}\NormalTok{, }\StringTok{\textquotesingle{}rewritten\textquotesingle{}}\NormalTok{)),}
\NormalTok{    engine TEXT }\KeywordTok{NOT} \KeywordTok{NULL}\NormalTok{,}
\NormalTok{    status TEXT }\KeywordTok{NOT} \KeywordTok{NULL} \KeywordTok{CHECK}\NormalTok{ (status }\KeywordTok{IN}\NormalTok{ (}\StringTok{\textquotesingle{}success\textquotesingle{}}\NormalTok{, }\StringTok{\textquotesingle{}error\textquotesingle{}}\NormalTok{, }\StringTok{\textquotesingle{}timeout\textquotesingle{}}\NormalTok{)),}
\NormalTok{    runtime\_ms }\DataTypeTok{INTEGER}\NormalTok{,}
\NormalTok{    rows\_returned }\DataTypeTok{INTEGER}\NormalTok{,}
\NormalTok{    result\_checksum TEXT,}
\NormalTok{    sample\_output TEXT,}
\NormalTok{    explain\_plan TEXT,}
\NormalTok{    estimated\_cost }\DataTypeTok{REAL}\NormalTok{,}
\NormalTok{    error\_message TEXT,}
\NormalTok{    executed\_at }\DataTypeTok{TIMESTAMP} \KeywordTok{DEFAULT} \FunctionTok{CURRENT\_TIMESTAMP}
\NormalTok{);}

\CommentTok{{-}{-} LLM analysis results}
\KeywordTok{CREATE} \KeywordTok{TABLE} \ControlFlowTok{IF} \KeywordTok{NOT} \KeywordTok{EXISTS}\NormalTok{ llm\_analyses (}
    \KeywordTok{id} \DataTypeTok{INTEGER} \KeywordTok{PRIMARY} \KeywordTok{KEY}\NormalTok{,}
\NormalTok{    query\_id }\DataTypeTok{INTEGER} \KeywordTok{NOT} \KeywordTok{NULL} \KeywordTok{REFERENCES}\NormalTok{ queries(}\KeywordTok{id}\NormalTok{),}
\NormalTok{    model TEXT }\KeywordTok{NOT} \KeywordTok{NULL}\NormalTok{,}
\NormalTok{    prompt\_version TEXT,}
\NormalTok{    score }\DataTypeTok{INTEGER} \KeywordTok{CHECK}\NormalTok{ (score }\OperatorTok{\textgreater{}=} \DecValTok{0} \KeywordTok{AND}\NormalTok{ score }\OperatorTok{\textless{}=} \DecValTok{100}\NormalTok{),}
\NormalTok{    score\_reason TEXT,}
\NormalTok{    issues\_found TEXT,}
\NormalTok{    input\_tokens }\DataTypeTok{INTEGER}\NormalTok{,}
\NormalTok{    output\_tokens }\DataTypeTok{INTEGER}\NormalTok{,}
\NormalTok{    llm\_cost\_usd }\DataTypeTok{REAL}\NormalTok{,}
\NormalTok{    latency\_ms }\DataTypeTok{INTEGER}\NormalTok{,}
\NormalTok{    error\_message TEXT,}
\NormalTok{    analyzed\_at }\DataTypeTok{TIMESTAMP} \KeywordTok{DEFAULT} \FunctionTok{CURRENT\_TIMESTAMP}
\NormalTok{);}

\CommentTok{{-}{-} LLM{-}generated recommendations}
\KeywordTok{CREATE} \KeywordTok{TABLE} \ControlFlowTok{IF} \KeywordTok{NOT} \KeywordTok{EXISTS}\NormalTok{ recommendations (}
    \KeywordTok{id} \DataTypeTok{INTEGER} \KeywordTok{PRIMARY} \KeywordTok{KEY}\NormalTok{,}
\NormalTok{    query\_id }\DataTypeTok{INTEGER} \KeywordTok{NOT} \KeywordTok{NULL} \KeywordTok{REFERENCES}\NormalTok{ queries(}\KeywordTok{id}\NormalTok{),}
\NormalTok{    llm\_analysis\_id }\DataTypeTok{INTEGER} \KeywordTok{NOT} \KeywordTok{NULL} \KeywordTok{REFERENCES}\NormalTok{ llm\_analyses(}\KeywordTok{id}\NormalTok{),}
\NormalTok{    recommended\_query TEXT }\KeywordTok{NOT} \KeywordTok{NULL}\NormalTok{,}
\NormalTok{    improvement\_reason TEXT,}
\NormalTok{    expected\_improvement\_category TEXT,}
\NormalTok{    improvement\_suggestion TEXT,}
\NormalTok{    created\_at }\DataTypeTok{TIMESTAMP} \KeywordTok{DEFAULT} \FunctionTok{CURRENT\_TIMESTAMP}
\NormalTok{);}

\CommentTok{{-}{-} Before/after cost comparisons}
\KeywordTok{CREATE} \KeywordTok{TABLE} \ControlFlowTok{IF} \KeywordTok{NOT} \KeywordTok{EXISTS}\NormalTok{ cost\_comparisons (}
    \KeywordTok{id} \DataTypeTok{INTEGER} \KeywordTok{PRIMARY} \KeywordTok{KEY}\NormalTok{,}
\NormalTok{    query\_id }\DataTypeTok{INTEGER} \KeywordTok{NOT} \KeywordTok{NULL} \KeywordTok{REFERENCES}\NormalTok{ queries(}\KeywordTok{id}\NormalTok{),}
\NormalTok{    recommendation\_id }\DataTypeTok{INTEGER} \KeywordTok{NOT} \KeywordTok{NULL} \KeywordTok{REFERENCES}\NormalTok{ recommendations(}\KeywordTok{id}\NormalTok{),}
\NormalTok{    original\_execution\_id }\DataTypeTok{INTEGER} \KeywordTok{REFERENCES}\NormalTok{ query\_executions(}\KeywordTok{id}\NormalTok{),}
\NormalTok{    rewritten\_execution\_id }\DataTypeTok{INTEGER} \KeywordTok{REFERENCES}\NormalTok{ query\_executions(}\KeywordTok{id}\NormalTok{),}
\NormalTok{    original\_runtime\_ms }\DataTypeTok{INTEGER}\NormalTok{,}
\NormalTok{    rewritten\_runtime\_ms }\DataTypeTok{INTEGER}\NormalTok{,}
\NormalTok{    runtime\_improvement\_pct }\DataTypeTok{REAL}\NormalTok{,}
\NormalTok{    original\_rows }\DataTypeTok{INTEGER}\NormalTok{,}
\NormalTok{    rewritten\_rows }\DataTypeTok{INTEGER}\NormalTok{,}
\NormalTok{    rows\_match }\DataTypeTok{BOOLEAN}\NormalTok{,}
\NormalTok{    checksum\_match }\DataTypeTok{BOOLEAN}\NormalTok{,}
\NormalTok{    semantic\_match }\DataTypeTok{BOOLEAN}\NormalTok{,}
\NormalTok{    validation\_status TEXT }\KeywordTok{CHECK}\NormalTok{ (validation\_status }\KeywordTok{IN}\NormalTok{ (}\StringTok{\textquotesingle{}match\textquotesingle{}}\NormalTok{, }\StringTok{\textquotesingle{}mismatch\textquotesingle{}}\NormalTok{, }\StringTok{\textquotesingle{}rewrite\_failed\textquotesingle{}}\NormalTok{, }\StringTok{\textquotesingle{}not\_executed\textquotesingle{}}\NormalTok{)),}
\NormalTok{    llm\_total\_cost\_usd }\DataTypeTok{REAL}\NormalTok{,}
\NormalTok{    net\_cost\_improvement TEXT,}
\NormalTok{    notes TEXT,}
\NormalTok{    created\_at }\DataTypeTok{TIMESTAMP} \KeywordTok{DEFAULT} \FunctionTok{CURRENT\_TIMESTAMP}
\NormalTok{);}
\end{Highlighting}
\end{Shaded}
